% frontmatter/nivel2/como-usar.tex -- instrucciones para el adulto.
\section*{Cómo usar este cuaderno}

Este es el segundo cuaderno de la serie. Está pensado para una niña o un
niño que ya lee con soltura frases completas y compuestas -- por
ejemplo, porque ha terminado \textit{Aprendo a leer} -- y que ahora da el
paso siguiente: entender oraciones largas, con varias partes unidas por
\textit{que}, \textit{porque}, \textit{aunque}, \textit{cuando},
\textit{si} o \textit{para que}, y fijarse en lo que el texto dice de
verdad, no en lo que parece que dice. Leer con lupa.

Como el primero, tiene una página para cada día laborable del año,
también en vacaciones, y cada página sigue siempre el mismo orden:

\begin{itemize}[leftmargin=6mm, itemsep=1mm]
\item \textbf{Fecha}, \textbf{Leído} y \textbf{Notas}, para el adulto.
\item El \textbf{texto del día}, en el recuadro verde: unas 55 palabras
  al principio del curso y unas 120 al final, en párrafos. El título del
  recuadro dice cómo leerlo (en voz alta, dos veces, con la voz de cada
  personaje, primero en silencio...).
\item Una \textbf{actividad de análisis}: para hacerla hay que volver al
  texto. La regla de oro es que \textbf{la respuesta siempre está en el
  texto} -- el de ese día o, si lo dice la instrucción, el de los días
  anteriores de la misma semana. Si no la encuentra, no se la digas:
  pregúntale «¿dónde lo dice?» y deja que vuelva a leer.
\end{itemize}

\textbf{Las actividades.} \textbf{Dibuja con lupa}: dibujar exactamente
lo que describe el texto; no importa que el dibujo sea bonito, sino que
sea exacto (cuántos, de qué color, dónde), y al terminar se repasa con
las preguntas de abajo, texto en mano. \textbf{Mapa}: seguir unas
indicaciones en una cuadrícula y marcar un sitio o un camino.
\textbf{Tabla de pistas}: un problema de lógica; se pone ✓ o
✗ en cada casilla hasta que solo queda una posibilidad.
\textbf{Caza los errores}: un resumen del texto con fallos que hay que
encontrar y corregir. \textbf{Verdadero o falso}: y, si es falso, cómo es
de verdad. \textbf{Mensaje secreto}: descifrar un mensaje con un código
que explica el texto. \textbf{Responde}, \textbf{Ordena},
\textbf{Relaciona}, \textbf{Ficha}, \textbf{Compara}, \textbf{Adivina},
\textbf{Viñetas} y \textbf{Crea}: preguntas de «¿por qué?» y «¿cómo
lo sabes?», sucesos que ordenar, datos que sacar de un texto
informativo, dos cosas que comparar, adivinanzas, una historia en tres
dibujos y un poco de escritura propia.

\textbf{Un caso cada semana.} Las páginas cuentan un año más de la vida
de \textbf{Lucía}, \textbf{Dani}, \textbf{Toby} y su familia, y del
\textbf{Club de la Lupa}, el club de detectives que Lucía funda con sus
amigos Sofía y Hugo. Casi todas las semanas hay un misterio: las pistas
aparecen de lunes a jueves, y el viernes toca resolverlo --
\textbf{Resuelve el caso}: rodear la respuesta y escribir qué pistas lo
demuestran. La solución se confirma el lunes siguiente, y está en la
\textbf{clave de respuestas} del final del cuaderno. Mejor no mirarla
antes de que lo intente: está ahí para comprobar, o para dar una pista
pequeña cuando se atasque. Cada trimestre tiene además un misterio
grande que se resuelve en su última semana.

\textbf{El ritmo.} Diez o quince minutos al día. Si durante una semana
entera lee y resuelve las páginas sin ayuda, puede hacer dos en un día;
si el texto se le hace largo, leedlo juntos en voz alta primero, un
párrafo cada uno, y que luego lo relea sola. Lo que importa es que
siga siendo un juego: los detectives no tienen prisa, tienen cuidado.

Al final de cada trimestre hay una medalla, y al final del cuaderno, un
diploma.
\newpage
